\section*{Capítulo \ref{ch:g3:stones-pebbles-sand} --- Pedras, seixos e areia}

\begin{solution}{exo:g3:stones-pebbles-sand:1}
O \omterm{def:g3:stones-pebbles-sand:rock}{substrato rochoso} é a \omterm{def:g3:stones-pebbles-sand:rock}{rocha} maciça que fica embaixo do solo em toda
parte. Nós o vemos onde não há solo para escondê-lo: nas encostas íngremes
das \omterm{def:g1:landscapes-around-us:landform}{montanhas}, nos penhascos à \omterm{def:g2:water-on-earth:ocean}{beira-mar}, ao longo das estradas abertas no
meio de uma \omterm{def:g1:landscapes-around-us:landform}{colina}.
\end{solution}

\begin{solution}{exo:g3:stones-pebbles-sand:2}
\omterm{def:g3:stones-pebbles-sand:sand}{Lama}, \omterm{def:g3:stones-pebbles-sand:sand}{areia}, \omterm{def:g3:stones-pebbles-sand:pebble}{seixo}, matacão.
\end{solution}

\begin{solution}{exo:g3:stones-pebbles-sand:3}
1 milímetro é menos que 2 milímetros: \omterm{def:g3:stones-pebbles-sand:sand}{areia}. 3 milímetros é mais que
2 milímetros: \omterm{def:g3:stones-pebbles-sand:pebble}{cascalho}.
\end{solution}

\begin{solution}{exo:g3:stones-pebbles-sand:4}
Naturais: o \omterm{def:g3:stones-pebbles-sand:pebble}{seixo}, o bloco de granito (cortado pelas pessoas, mas a \omterm{def:g3:stones-pebbles-sand:rock}{rocha} é
natural), a \omterm{def:g3:stones-pebbles-sand:sand}{areia} da praia. Feitos pelas pessoas: o tijolo e o vidro.
\end{solution}

\begin{solution}{exo:g3:stones-pebbles-sand:5}
A água do \omterm{def:g1:landscapes-around-us:water}{rio} o fez rolar pelo fundo e bater em outras pedras durante
muito tempo: os cantos dele se gastaram.
\end{solution}

\begin{solution}{exo:g3:stones-pebbles-sand:6}
O granito, para um muro; lajes de pedra finas e planas (de uma \omterm{def:g3:stones-pebbles-sand:rock}{rocha} como a ardósia), para telhados.
\end{solution}

\begin{solution}{exo:g3:stones-pebbles-sand:7}
A redonda: quanto mais tempo uma pedra rola na água, mais os cantos dela se
gastam. A pontuda não viajou para longe da \omterm{def:g3:stones-pebbles-sand:rock}{rocha} de onde se soltou.
\end{solution}

\begin{solution}{exo:g3:stones-pebbles-sand:8}
O maior \omterm{def:g3:stones-pebbles-sand:pebble}{seixo} tem 3 centímetros de comprimento. Um centímetro tem 10
milímetros, e $10 \div 2 = 5$: cinco grãos de 2 milímetros cabem ao longo
de 1 centímetro.
\end{solution}

\begin{solution}{exo:g3:stones-pebbles-sand:9}
Na lagoa calma. A água rápida mantém os grãos minúsculos de \omterm{def:g3:stones-pebbles-sand:sand}{lama} em
movimento; só onde a água está quase parada eles conseguem se depositar no fundo.
\end{solution}

\begin{solution}{exo:g3:stones-pebbles-sand:10}
Porque eles são feitos de materiais tirados do chão: os tijolos, de argila;
o concreto, de \omterm{def:g3:stones-pebbles-sand:sand}{areia}, \omterm{def:g3:stones-pebbles-sand:pebble}{cascalho} e cimento; o vidro, de \omterm{def:g3:stones-pebbles-sand:sand}{areia}. A argila, a
\omterm{def:g3:stones-pebbles-sand:sand}{areia} e o \omterm{def:g3:stones-pebbles-sand:pebble}{cascalho} vêm de \omterm{def:g3:stones-pebbles-sand:rock}{rocha} quebrada.
\end{solution}
